\section{The discrete moments for a fixed primitive character}
\label{sec:dirichlet-moments}
This section proves the contour formulas used in
Proposition~\ref{prop:square-moments}. Keeping the character phases in
the additive transform is essential: their cancellation, followed by
character orthogonality, gives the same leading constants for every
fixed primitive character.

Throughout, $q$ is fixed, $Z_q(s)=\zeta(s)\prod_{p\mid q}(1-p^{-s})$,
$E_q=\varphi(q)/q$, and coefficient polynomials have the form
$A_\chi(s)=\sum_{n\le y}a(n)\chi(n)n^{-s}$, where the real coefficients
satisfy $|a(n)|\le C\tau_r(n)(\log(2n))^r$ for fixed $C,r$.
Their supports satisfy $y\le T^\vartheta$ for a fixed $\vartheta<1/2$.
In particular, the squared mollifier has
$|a_T(n)|\le C_P^2\tau_2(n)$. In the character averages below its
subsequence $a_T(eh)$ is kept exactly and bounded by
$C_P^2\tau_2(e)\tau_2(h)$. This follows from the divisor envelope
and $\tau_2(eh)\le\tau_2(e)\tau_2(h)$; the argument requires no
product inequality for the coefficients $a_T$ themselves.
All shift variables have modulus at most a fixed constant divided by
$\log T$; all assertions below are uniform on slightly larger shift
circles, so Cauchy differentiation is permitted.

\subsection{Bounds for auxiliary characters}
We first record explicitly the consequence of the zero-location input
used for the growing auxiliary moduli. For every fixed $\eta,\epsilon>0$,
\begin{equation}\label{eq:uniform-character-growth}
 |L(\sigma+it,\psi)|+|L(\sigma+it,\psi)^{-1}|
 \ll_{\eta,\epsilon}\{Q(|t|+3)\}^{\epsilon},
 \qquad \tfrac12+\eta\le\sigma\le2,
\end{equation}
for nonprincipal characters modulo $Q$. For principal characters the
same bound holds on the line $\sigma=1/2+\eta<1$, and on the horizontal
parts of the strip when $|t|\ge2$.

Here is a proof including uniformity in $Q$. For primitive nonprincipal
$\psi$, the real part of the completed logarithmic derivative at $2+it$
is the positive zero sum
\[
 \sum_\rho\frac{3/2}{(3/2)^2+(t-\gamma)^2}
     =\Re\frac{\Lambda'}\Lambda(2+it,\psi)
     \ll\log\mathcal C,\qquad \mathcal C=Q(|t|+3).
\]
In the Hadamard product centered at $1/2$, the linear coefficient is
purely imaginary by conjugate reflection, and the centered zero
reciprocals have zero real part; this proves the equality. The Euler
series at two and the gamma logarithmic derivative prove the bound.
Subtracting logarithmic derivatives at $\sigma+it$ and $2+it$ then gives
$|L'/L(\sigma+it,\psi)|\ll_\eta\log\mathcal C$.
For $x>1$, contour integration of $-(L'/L)(s+w)x^w/w^2$ gives
\begin{align*}
 -\frac{L'}L(s,\psi)={}&
 \sum_{n\le x}\frac{\Lambda(n)\psi(n)}{n^s}
                     \frac{\log(x/n)}{\log x}
 +\frac1{\log x}\left(\frac{L'}L\right)'(s,\psi)\\
 &+\frac1{\log x}\sum_\rho\frac{x^{\rho-s}}{(\rho-s)^2}
 +\frac1{\log x}\sum_{\rho_{\rm triv}}
                       \frac{x^{\rho_{\rm triv}-s}}{(\rho_{\rm triv}-s)^2}.
\end{align*}
The residues at zeros are counted with multiplicity. The functional
equation makes a left edge tend to zero; horizontal edges can be
chosen away from ordinates using the positive zero sum, which gives
$L'/L=O(\log^2(QH))$ there. The quadratic denominators make both zero
sums absolutely convergent.
Integrate this identity in the real coordinate from $\sigma$ to infinity,
and choose $x=\max(3,\log\mathcal C)$. The prime term is $O(\sqrt x)$;
the derivative term is $O_\eta(\log\mathcal C/\log x)$; the nontrivial
zero term is
$O_\eta(x^{-\eta}\log\mathcal C/(\log x)^2)$; the trivial-zero term
is bounded. It follows that
\[
 |\log L(s,\psi)|\ll_\eta\sqrt{\log\mathcal C}
       +\frac{\log\mathcal C}{\log\log(\mathcal C+3)}
       =o_\eta(\log\mathcal C).
\]
Exponentiation proves \eqref{eq:uniform-character-growth}. Passing to
an imprimitive character adds deleted Euler factors and their inverses,
whose product is $O_{\eta,\epsilon}(Q^\epsilon)$: absorb each sufficiently
large prime factor into $p^\epsilon$ and bound the finite remaining
product separately. For the principal character the pole at one adds
a bounded or $O(\sqrt x)$ term in the indicated regions. Fixed-order
derivatives follow from Cauchy's formula on a smaller strip.

For a fixed character these estimates also imply the critical-line
bound $L_\chi^{(j)}(1/2+it)\ll_{\chi,j,\epsilon}(|t|+3)^\epsilon$.
To see this, use \eqref{eq:uniform-character-growth} on
$\Re s=1/2+\eta$, reflect it to $1/2-\eta$ with the functional
equation, and apply the three-lines theorem in that strip. Its
reflection costs at most $T^{\eta+\epsilon}$; choose $\eta$ smaller
than the desired exponent, and apply Cauchy on a fixed smaller
strip. This also supplies the bounds on the horizontal contour edges
used below.

\subsection{Averaging the additive transform}
Write $e(u)=e^{2\pi iu}$. The stationary transform of the functional
equation leads to the following expression, for $T\le U\le2T$:
\begin{equation}\label{eq:additive-transform}
 \frac{\tau(\bar\chi)}q
 \sum_{\substack{k\le y\\(k,q)=1}}
   \frac{a_2(k)\bar\chi(k)}{k^{1-\beta}}
 \sum_{n\le qUk/(2\pi)}\frac{\chi(n)b(n)}{n^\beta}
                  e\left(-\frac n{qk}\right).
\end{equation}
Here $b=c*a_1$, where $c$ is a convolution of at most four shifted
constant or M\"obius sequences, with at most one M\"obius factor.
Fixed-order shift derivatives are allowed.

\begin{lemma}\label{lem:dirichlet-additive}
All nonprincipal character components of
\eqref{eq:additive-transform} have total
$O_{q,\epsilon}(T^{1/2+\epsilon}y)$. Its principal component is
exactly
\begin{equation}\label{eq:principal-transform}
 \frac1{\varphi(q)}
 \sum_{\substack{k\le y\\(k,q)=1}}\frac{a_2(k)}{k^{1-\beta}}
 \sum_{\substack{n\le qUk/(2\pi)\\(n,q)=1}}
     \frac{b(n)}{n^\beta}
       \frac{\mu(k/(n,k))}{\varphi(k/(n,k))}.
\end{equation}
Replacing its inner sums by the residues to the right of
$\Re s=1/2+\eta$, for a fixed sufficiently small $\eta>0$, has
the same total error bound.
\end{lemma}

\begin{proof}
Put $d=(n,k)$, $k'=k/d$, $n=dm$ and $Q=qk'$. On the units modulo
$Q$, expand
\[
 \chi(m)e(-m/Q)=\sum_{\psi\bmod Q}u_\psi\psi(m),
 \qquad \sum_\psi|u_\psi|^2=1.
\]
Chinese remaindering and the primitive Gauss sum give
\[
 u_{\psi_0}=\frac{\mu(k')\chi(k')\tau(\chi)\bar\chi(-1)}
                      {\varphi(q)\varphi(k')}.
\]
Using $\tau(\chi)\tau(\bar\chi)=\chi(-1)q$ and
$\chi(d)\chi(k')=\chi(k)$ cancels every character phase in the
principal scalar, leaving $\mu(k')/[\varphi(q)\varphi(k')]$.
This proves \eqref{eq:principal-transform}.

For the error, factor the relevant Dirichlet series exactly as
\begin{align*}
 \sum_{(m,Q)=1}\frac{b(dm)\psi(m)}{m^s}
 =\sum_{e\mid d}&\left(
 \sum_{\substack{h\le y/e\\(h,Qd/e)=1}}
              \frac{a_1(eh)\psi(h)}{h^s}\right)\left(
 \sum_{(m,Q)=1}\frac{c((d/e)m)\psi(m)}{m^s}\right).
\end{align*}
For a divisor $h_0\mid dm$, this is the bijection
$e=(h_0,d)$, $h_0=eh$, $(h,d/e)=1$, followed by $m=hm'$.
The second factor is a product of shifted $L$-functions and one
possible reciprocal, times local modifications at primes dividing
$d/e$. Before a valuation shift, the local series is
$C_p(z)=A_p(z)/\prod_{j=1}^{d_p}(1-b_jz)$, where $d_p\le4$,
$A_p$ is either $1$ or $1-a z$, and in the latter case $d_p\le3$.
All $a,b_j$ have modulus $p^{O(1/\log T)}$.
The case consisting of a single shifted M\"obius factor is handled
separately: $C_p(z)=1-a z$, its valuation-one shifted series is
$-a$, and every higher shifted series is zero. Its nonzero ratio
to $C_p$ is $-a/(1-a z)$ and satisfies the bound below. An empty
convolution has zero shifted series. In every remaining case
$d_p\ge1$.
For $v\ge1$, multiply the shifted series
$\sum_{n\ge0}c(p^{n+v})z^n$ by the same denominator.
The coefficient recurrence for $C_p$ makes every coefficient of
degree at least $d_p$ vanish. The remaining numerator therefore has
degree at most $d_p-1$; it need not be linear.
Its coefficients are bounded by
$O((v+1)^3p^{O((v+1)/\log T)})$, by expanding the at most four
geometric factors. Dividing by $C_p$ adds only the possible factor
$(1-a z)^{-1}$. Thus shifting by $v$ costs
$O((v+1)^3p^{O(v/\log T)})$; on
$\Re s\ge1/2+\eta$, every added denominator stays away from zero.
Thus \eqref{eq:uniform-character-growth} bounds its multiplier by
$\tau_4(d/e)\{dQ(|\Im s|+3)\}^\epsilon$.

Orthogonality and Cauchy in residue classes bound the finite polynomial
by
\begin{align*}
 \sum_{\psi\bmod Q}\left|
 \sum_{\substack{h\le y/e\\(h,Qd/e)=1}}
       a_1(eh)\psi(h)h^{-s}\right|^2
 &\le\varphi(Q)\left(1+\frac y{eQ}\right)
                      \sum_{h\le y/e}|a_1(eh)|^2h^{-2\Re s}\\
 &\ll_{\eta,r}\tau_r(e)^2(\log(2y))^{2r}(Q+y/e).
\end{align*}
Here $\tau_r(eh)\le\tau_r(e)\tau_r(h)$ and
$\sum\tau_r(h)^2h^{-1-2\eta}<\infty$. Together with
$\sum|u_\psi|^2=1$, this bounds the combined integrand by
\[
 \{dQ(|\Im s|+3)T\}^\epsilon\tau_{r+4}(d)\sqrt{Q+y}.
\]
Use truncated Perron inversion at $x=qUk/(2\pi d)$, replacing the
cutoff temporarily by $\lfloor x\rfloor+1/2$, with height $H=T^{10}$.
The initial line is $1+C/\log T$, beyond all shifts. The elementary
kernel is
\[
 \frac1{2\pi i}\int_{c-iH}^{c+iH}\frac{u^s}{s}ds
 =\mathbf1_{u>1}+O\left(u^c\min\{1,(H|\log u|)^{-1}\}\right).
\]
It follows by closing the infinite vertical contour on the appropriate
side and integrating its tails by parts. Half-integrality and the
divisor coefficient bound make the near-cutoff error
$O(d^\epsilon T^\epsilon x/H)$; the remaining absolute Dirichlet
norm is a fixed power of $\log T$ and is harmless. Apply the kernel
after combining characters: the original multiplier has modulus one,
and its nonprincipal remainder has modulus at most two, so no number
of characters enters the truncation error.

Shift to $1/2+\eta$. Nonprincipal characters contribute no poles;
record the explicit shifted poles for the principal one. The preceding
integrand bound, and $x^s/s$, give
\[
 O_{q,\epsilon}\left(T^\epsilon(Tk/d)^{1/2}
                       \tau_{r+4}(d)\sqrt{qk/d+y}\right).
\]
Choose the preliminary exponents and $\eta$ small enough to absorb
$H^\epsilon$, $x^\eta$ and logarithms. Horizontal edges have the
factor $H^{-1}$ and satisfy the same bound. The residues change by
at most $T^\epsilon d^\epsilon$ times fixed divisor factors when the
cutoff is restored. Multiplying by the outer $1/k$ and summing gives
\[
 T^{1/2+\epsilon}\sqrt y\sum_{k\le y}k^{-1/2}|a_2(k)|
                         \sum_{d\mid k}d^{-1/2}\tau_{r+4}(d)
 \ll T^{1/2+\epsilon}y.
\]
Slightly enlarging the shift circles allows differentiation; endpoint
partial summation allows fixed logarithmic weights.
\end{proof}

\subsection{The three residues and the discrete second moment}
Choose $c=1+C_0/L$, with the fixed $C_0$ larger than all scaled
shift bounds, and a fixed lower ordinate $t_0>2\pi$. Define
\[
 \mathcal K(\alpha,\beta,\gamma;U)
 =\frac1{2\pi i}\int_{c+it_0}^{c+iU}
 \frac{L_\chi(s+\gamma)}{L_\chi(s)}L_\chi(s+\alpha)
 L_{\bar\chi}(1-s+\beta)A_\chi(s)A_{\bar\chi}(1-s)\,ds.
\]
Applying $L_{\bar\chi}(1-s+\beta)=
\Delta_\chi(s-\beta)^{-1}L_\chi(s-\beta)$ and expanding absolutely
on the right line gives \eqref{eq:additive-transform}, with
$c=\mu*n^{-\gamma}*n^{-\alpha}*n^\beta$ in its coefficient notation,
and error $O_{q,\epsilon}(T^{1/2+\epsilon}y)$.
For complex $\beta$, make the change of variable $z=s-\beta$ before
applying the stationary-phase formula. For each pair of coefficient
indices,
\[
 (m/k)^{-s}\Delta_\chi(s-\beta)^{-1}
   =(m/k)^{-\beta}(m/k)^{-z}\Delta_\chi(z)^{-1}.
\]
The new vertical line has real part $c-\Re\beta$, which belongs
to a fixed compact interval about one and remains to the right of
every Dirichlet-series convergence boundary. Thus the amplitude
exponent in the cited formula is the real number
$c-\Re\beta-1/2$. Restoring the endpoint ordinates from
$t_0-\Im\beta,U-\Im\beta$ to $t_0,U$ adds segments of length
$O(L^{-1})$. The absolute Dirichlet norm and Stirling's bound give
total contribution $O(\sqrt{qT}\,yL^{C_1})$, with the bottom
contributing only $O(\sqrt q\,yL^{C_1})$, uniformly for $q\le T$.
Both lie within the stated errors. The factor $(m/k)^{-\beta}$
gives exactly $m^{-\beta}k^\beta$ in
\eqref{eq:additive-transform}.

To verify the error explicitly, the stationary point for inner index
$m$ and outer index $k$ is $t=2\pi m/(qk)$; the gamma quotient
contributes scalar $\tau(\bar\chi)/q$ and phase $e(-m/(qk))$.
The $O(1/t)$ relative Stirling error is $O_q(y\sqrt T\log^C T)$.
Outside a factor-two neighborhood of the stationary point, integration
by parts has the same bound. Inside it, the stationary formula stated
in the proof of Lemma~\ref{lem:mixed-phase} is applied only when
$m/(qk)>1$. The remaining terms are treated on a bounded lower
interval and by the first-derivative test above it, with the same
total error. For $m/(qk)>1$ the stationary formula gives errors bounded,
for each outer coefficient, by
\[
 C_q k^{-1/2}\sum_{m\le C_qkT}\frac{\tau_R(m)}{\sqrt m}
 \ll\sqrt T(kT)^\epsilon
\]
and
\[
 C_q\sqrt T(kT)^\epsilon\sum_{m\asymp qkT}
       (|m-qkU/(2\pi)|+qk\sqrt T)^{-1}
 \ll\sqrt T(kT)^\epsilon\log T.
\]
Connecting the shifted contour endpoints costs the same amount.
Summing the divisor-bounded outer coefficients proves the assertion.

For coprime $h,k$ put $f_{\alpha,\gamma}=\mu*n^{-\alpha}*n^{-\gamma}$
and
\[
 Z^{(q)}_{\alpha,\beta,\gamma,h,k}
 =h^{-\beta}\frac{Z_q(1+\alpha+\beta)Z_q(1+\beta+\gamma)}{Z_q(1+\beta)}
 \prod_{p^v\parallel k}
 \frac{\sum_{m\ge0}f_{\alpha,\gamma}(p^{m+v})p^{-m(1+\beta)}}
      {\sum_{m\ge0}f_{\alpha,\gamma}(p^m)p^{-m(1+\beta)}}.
\]
Then $\mathcal K$ equals
\begin{equation}\label{eq:dirichlet-three-residue}
 \frac1{2\pi}\int_0^U\sum_{\substack{(h,k)=1\\(ghk,q)=1}}
 \frac{a(gh)a(gk)}{ghk}
 \left\{\begin{aligned}
 &Z^{(q)}_{\alpha,\beta,\gamma,h,k}\\
 &+(qt/(2\pi))^{-\alpha-\beta}Z^{(q)}_{-\beta,-\alpha,\gamma,h,k}\\
 &+(qt/(2\pi))^{-\beta-\gamma}Z^{(q)}_{\alpha,-\gamma,-\beta,h,k}
 \end{aligned}\right\}dt
 +O_{q,\epsilon}(T^{1/2+\epsilon}y).
\end{equation}
Here is the residue computation. In the principal transform unfold
$b=c*a$, and express the two coefficient indices as $gh,gk$ with
$(h,k)=1$. The remaining Dirichlet series is
\[
 \sum_{(n,q)=1}c(n)n^{-s-\beta}
                \frac{\mu(k/(n,k))}{\varphi(k/(n,k))}.
\]
Its unrestricted product is
$Z_q(s)Z_q(s+\beta+\gamma)Z_q(s+\beta+\alpha)/Z_q(s+\beta)$.
At $p^v\parallel k$ its weight is zero for exponents below $v-1$,
$-1/(p-1)$ at $v-1$, and one thereafter. There are precisely the
three possible poles $1$, $1-\alpha-\beta$, $1-\beta-\gamma$.
For the first, let $F_p(X)=\sum f_{\alpha,\gamma}(p^j)X^j$,
$C_p(X)=F_p(X)/(1-p^\beta X)=\sum c_jX^j$, and $X=p^{-1-\beta}$.
The recurrence $c_j=f_{\alpha,\gamma}(p^j)+p^\beta c_{j-1}$ gives
\[
 (1-1/p)\left(\sum_{j\ge v}c_jX^j-
                         \frac{c_{v-1}X^{v-1}}{p-1}\right)
             =\sum_{j\ge v}f_{\alpha,\gamma}(p^j)X^j.
\]
This proves the displayed local factor after restoring the cutoff.
The residue $E_q$ cancels $q/\varphi(q)$ exactly.
The substitutions $(\alpha,\beta,\gamma)\mapsto
(-\beta,-\alpha,\gamma)$ and $(\alpha,-\gamma,-\beta)$ move each
other pole to the first and yield the remaining terms by
$x^{1-z}/(1-z)=\int_0^x t^{-z}dt$. Coinciding-shift values are
removable, since the original finite-coefficient integral is holomorphic.

For use in \eqref{eq:square-euler-F}--\eqref{eq:square-euler-G}, we
state the resulting discrete moment explicitly. Put
\begin{align*}
 \mathcal J_q(\alpha,\beta;U)
 =\frac1{2\pi}\int_0^U&\log\frac{qt}{2\pi}
 \sum_{\substack{(h,k)=1\\(ghk,q)=1}}\frac{a(gh)a(gk)}{ghk}\\
 &\cdot\left\{\frac{Z_q(1+\alpha+\beta)}{h^\beta k^\alpha}
 +(qt/(2\pi))^{-\alpha-\beta}
              \frac{Z_q(1-\alpha-\beta)}{h^{-\alpha}k^{-\beta}}\right\}dt.
\end{align*}
Let $\mathcal I_q(\alpha,\beta;U)$ be the $\gamma$ derivative at zero
of the main term in \eqref{eq:dirichlet-three-residue}. Then
\begin{align}\label{eq:dirichlet-discrete}
 \sum_{T<\gamma\le2T}&L_\chi(\rho+\alpha)
 L_{\bar\chi}(1-\rho+\beta)A_\chi(\rho)A_{\bar\chi}(1-\rho)\nonumber\\
 &=\bigl[\mathcal J_q(\alpha,\beta;U)+\mathcal I_q(\alpha,\beta;U)
        +\mathcal I_q(\beta,\alpha;U)\bigr]_{U=T}^{2T}
       +O_{q,\epsilon}(T^{1/2+\epsilon}y).
\end{align}
Indeed, integrate $(L_\chi'/L_\chi)(s)$ times the four displayed
factors over a rectangle. The right edge is
$\partial_\gamma\mathcal K|_{\gamma=0}$.
On the left use $L_\chi'/L_\chi=\Delta_\chi'/\Delta_\chi-
L_{\bar\chi}'/L_{\bar\chi}(1-s)$ and reflect. Conjugation and
reflection transpose the shifts; the principal expression in
\eqref{eq:dirichlet-three-residue} has real coefficients, which yields
$\mathcal I_q(\beta,\alpha)$. The gamma logarithmic derivative is
$-\log(qt/(2\pi))+O_q(1/t)$. Its main integral can be evaluated by
setting $\gamma=0$ in \eqref{eq:dirichlet-three-residue}; the quotient
$L(s+\gamma)/L(s)$ becomes one, and the third term vanishes through
$1/Z_q(1)$. Endpoint partial summation supplies the logarithm and gives
$\mathcal J_q$.

For completeness the horizontal edges satisfy
$L_\chi'/L_\chi=O_q(\log^2T)$ at heights avoiding zeros, and
\[
 |L_\chi(s+\alpha)L_{\bar\chi}(1-s+\beta)|
       \ll_{q,\epsilon}T^{|\Re s-1/2|+\epsilon},\qquad
 |A_\chi(s)A_{\bar\chi}(1-s)|\ll y^{1+\epsilon}
\]
through the strip. Their contribution is
$O_{q,\epsilon}(yT^{1/2+\epsilon})$.
Restoring the endpoints $T,2T$ costs
$O_{q,\epsilon}(yT^\epsilon\log T)$, using the local zero count and
the critical-line bound. All residues count multiplicity. Thus the
formula and its differentiated versions apply even when there are
multiple zeros. Taking $A=B_{\chi,T}^2$ and differentiating both
shifts proves the moment interface used in the kernel calculation.

\subsection{The first moment and the signed mean}
To prove \eqref{eq:square-first}, integrate
$L_\chi'(s)^2A_\chi(s)/L_\chi(s)$ with
$A_\chi=B_{\chi,T}^2$. At a multiple zero the integrand is holomorphic
and the target summand vanishes. This pole calculation for
$L_\chi'^2/L_\chi$ is also the simple-zero transform in
Banks \cite[Section~2, p.~3]{banks-simple-transform}; the polynomial
factor here inserts the specified mollifier. The right-edge series has no constant
term and contributes $O(y^{1+\epsilon}\log^C T)$. Reflection of the
left edge, followed by the stationary transform, gives a negative
principal transform with inner coefficient
\[
 c_k(n)=(\log k)^2-(\log n)^2+(\Lambda*\log)(n).
\]
This follows directly from
\[
 \Delta_\chi(s)^{-1}\{D_\chi(s)^2L_\chi(s)
       -2D_\chi(s)L_\chi'(s)+L_\chi'(s)^2/L_\chi(s)\},
 \qquad D_\chi=\Delta_\chi'/\Delta_\chi,
\]
after replacing $D_\chi$ by $-\log(qt/(2\pi))$ and evaluating the
logarithm at the stationary point. The reflected expression is
conjugated; its principal residues are real, so their values are
unchanged. All remaining components have the error in
Lemma~\ref{lem:dirichlet-additive}.

The principal series is
\[
 W_k(s)=\sum_{(n,q)=1}c_k(n)n^{-s}
              \frac{\mu(k/(n,k))}{\varphi(k/(n,k))}.
\]
For its residues write
\[
 F_k(s)=Z_q(s)A_k(s),\quad
 A_k(s)=\prod_{p^v\parallel k}
             \frac{p^{-(v-1)s}(p^{1-s}-1)}{p-1}.
\]
The coefficient identity
$(\Lambda*\log)(n)-(\log n)^2=-\sum_{p^j\mid n}j(\log p)^2$
is obtained from $Z_q'^2/Z_q-Z_q''=-Z_q(Z_q'/Z_q)'$.
Locally its insertion multiplies the coefficient of $X^j$ by
$j(j+1)/2$, that is, applies $(X^2\partial_X^2+2X\partial_X)/2$.
Let $z=s-1$, $Z_q(s)=E_q/z+C_q+O_q(z)$. Direct expansion gives
\begin{align*}
 W_1(s)&=-E_qz^{-3}+O_q(z^{-2}),\\
 W_p(s)&=\frac{E_q\log p}{p-1}z^{-2}
 +\left\{-E_q\left(\frac{3\log^2p}{2(p-1)}
        +\frac{\log^2p}{(p-1)^2}\right)
        +\frac{C_q\log p}{p-1}\right\}z^{-1}+O_{p,q}(1),\\
 W_{pr}(s)&=-\frac{E_q\log p\log r}{(p-1)(r-1)}z^{-1}+O_{p,r,q}(1)
 \quad(p\ne r).
\end{align*}
These are obtained by applying the displayed differential operator
to the rational local geometric factor. If $\omega(k)\ge3$, its
$\omega(k)$ zeros at one cancel every possible pole. Terms involving
repeated prime powers contribute $O_{q,P}(L)$ after division by
$U/(2\pi)$: their residues are bounded by
\[
 C|a(p^v)|p^{-v}(Lv^2\log p+v^3\log^2p),\quad v\ge2,
\]
and by $C|a(p^vr^w)|p^{-v}r^{-w}v^2w^2\log p\log r$ for two
primes. A repeated-prime sum converges absolutely; the other costs
at most $O(L)$. Replacing $p-1$ by $p$ also costs $O(L)$.
Taking residues of $W_k(s)x^s/s$ with $x=qUk/(2\pi)$ and its
negative outer scalar yields
\[
 \frac U{2\pi}\left\{\tfrac12\log^2U
 +\sum_p\frac{a(p)}p(-\log U\log p+\tfrac12\log^2p)
 +\sum_{p<r}\frac{a(pr)}{pr}\log p\log r+O_{q,P}(L)\right\}.
\]
The conductor cancels because $qE_q/\varphi(q)=1$; deleting its
finitely many primes affects only the error. This proves the first
residue expression used in Proposition~\ref{prop:square-moments}.

We finally justify the signed cancellation. For two polynomials with
real underlying coefficients and lengths at most $y<T^{1/2}$,
\begin{equation}\label{eq:real-bilinear}
 \Im\sum_{T<\gamma\le2T}L_\chi'(\rho)
             X_\chi(\rho)Y_{\bar\chi}(1-\rho)
       =O_{q,\epsilon}(T^{1/2+\epsilon}y).
\end{equation}
Integrate $L_\chi'^2/L_\chi$ times these polynomials on the same
rectangle. On the right, the diagonal relation $nh=k$ has character
factor $\chi(n)\chi(h)\bar\chi(k)=1$ and real coefficient.
The off-diagonal integral is $O(y^{1+\epsilon}\log^C T)$: group
$m=nh$, split $m$ outside or inside $[k/2,2k]$, and in the latter
range use $|\log(m/k)|\gg|m-k|/k$ and its harmonic distance sum.
The reflected left edge is the three-term expression just displayed,
with the outer and inner polynomials exchanged. Its stationary
principal component is real by the exact scalar
\eqref{eq:principal-transform}; at the stationary point each gamma
logarithm is the real number $\log(m/k)$. The nonprincipal components,
horizontal edges and restored endpoints have the error above. This
proves \eqref{eq:real-bilinear}, retaining complex characters throughout.

Choose a fixed $\delta>0$ with $\nu+\delta<1/2$, and apply
\eqref{eq:real-bilinear} to
\[
 X_\chi=B_{\chi,T},\qquad
 Y_{\bar\chi}(s)=B_{\bar\chi,T}(s)
            \sum_{m\le T^\delta}\bar\chi(m)m^{-1-s}.
\]
The exterior tail is at most $2T^{-\delta/2}$. By the second moment,
its total contribution is bounded by
\[
 2T^{-\delta/2}
 \left(N_{\chi,T}\sum_{T<\gamma\le2T}
             |L_\chi'(\rho)B_{\chi,T}(\rho)^2|^2\right)^{1/2}
       =o(N_{\chi,T}L).
\]
If $L_\chi'(\rho)\overline{L_\chi(\rho+1)}=u+iv$, then exactly
$K_{\rho,\chi}=v+u/(2\gamma)$. The imaginary part is controlled
by \eqref{eq:real-bilinear}; the last term has weighted sum
$O(N_{\chi,T}L/T)$ by Cauchy and the Euler product bound.
This proves \eqref{eq:square-cancel}, and completes the proof of
Proposition~\ref{prop:square-moments}.
